\section*{الفصل \ref{ch:g10:synthesis-yield} --- التخليق: المردود والنقاوة}

\begin{solution}{exo:g10:synthesis-yield:1}
$\eta = 4.2 / 5.0 = 0.84$، أي \omterm{def:g10:synthesis-yield:yield}{مردود} قدره \qty{84}{\percent}.
\end{solution}

\begin{solution}{exo:g10:synthesis-yield:2}
وزن \omterm{def:g7:chemical-reactions:reactant}{المتفاعلات}؛ \omterm{def:g10:synthesis-yield:reflux}{والتسخين بالارتداد}؛ \omterm{def:g10:synthesis-yield:vacuum-filtration}{والترشيح تحت التفريغ} \omterm{def:g10:synthesis-yield:synthesis}{للناتج الخام}؛ \omterm{def:g10:synthesis-yield:recrystallisation}{وإعادة
التبلور}؛ وقياس درجة الانصهار.
\end{solution}

\begin{solution}{exo:g10:synthesis-yield:3}
إذا دخل الماء من الأسفل ملأ الغلاف كله قبل أن يخرج من الأعلى، فيُبرَّد المكثّف
على طوله كله. ويُترك الأعلى مفتوحًا حتى لا يكون الجهاز مغلقًا: فتسخين جهاز
مغلق يرفع الضغط فيه وقد يفجّره.
\end{solution}

\begin{solution}{exo:g10:synthesis-yield:4}
يجب أن \omterm{def:g2:mixing-and-dissolving:dissolve}{يذوب} \omterm{def:g7:chemical-reactions:reactant}{الناتج} كثيرًا في \omterm{def:g6:solutions-and-solubility:solute}{المذيب} الساخن ولا يكاد \omterm{def:g2:mixing-and-dissolving:dissolve}{يذوب} في \omterm{def:g6:solutions-and-solubility:solute}{المذيب} البارد.
والشوائب، الموجودة بمقادير صغيرة، تبقى ذائبة في \omterm{def:g6:solutions-and-solubility:solute}{المذيب} البارد وتنتهي في
\omterm{def:g3:separating-mixtures:filter}{الراشح}.
\end{solution}

\begin{solution}{exo:g10:synthesis-yield:5}
لا: فهي تنصهر تحت \qty{135}{\celsius} وعلى مدى عدة درجات، وهذه علامة جسم صلب
غير نقي. فيجب أن يُعاد تبلورها، وتُجفف، ويُعاد قياس درجة
انصهارها.
\end{solution}

\begin{solution}{exo:g10:synthesis-yield:6}
$n = 2.00 / 138.0 = \qty{0.0145}{mol}$ من حمض الساليسيليك، وهو \omterm{def:g10:reaction-progress-table:limiting}{المتفاعل
المحدِّد}؛ $m_{\max} = 0.0145 \times 180.0 = \qty{2.61}{g}$؛
$\eta = 2.03 / 2.61 = 0.78$، أي \qty{78}{\percent}.
\end{solution}

\begin{solution}{exo:g10:synthesis-yield:7}
الكتلة الموزونة تشمل الماء، فهي أكبر من كتلة الأسبرين: فيكون \omterm{def:g10:synthesis-yield:yield}{المردود} المحسوب
أعلى من اللازم، بل يتجاوز هنا \qty{100}{\percent} على نحو مستحيل. \omterm{def:g10:synthesis-yield:synthesis}{والناتج
الخام} الجاف لا يزال يحوي شوائب (حمض ساليسيليك لم يتفاعل، ونواتج ثانوية) تضيف
هي أيضًا كتلة: فيكون مردوده أيضًا أعلى من \omterm{def:g10:synthesis-yield:yield}{المردود} الحقيقي للأسبرين
النقي.
\end{solution}

\begin{solution}{exo:g10:synthesis-yield:8}
هو أسرع بكثير، ويترك البلورات أجف بكثير، لأن الامتصاص يسحب معظم السائل؛
كما يسهل غسل البلورات على القمع وكشطها عن ورق \omterm{def:g3:separating-mixtures:filter}{الترشيح}
المسطح.
\end{solution}

\begin{solution}{exo:g10:synthesis-yield:9}
\omterm{def:g7:chemical-reactions:reactant}{الناتج} $3.1/5.0 = 0.62$؛ والأسبرين $3.1/5.0 = 0.62$؛ وحمض الساليسيليك
$2.2/5.0 = 0.44$. \omterm{def:g7:chemical-reactions:reactant}{فالناتج} يتنقل كالأسبرين ولا يُظهر بقعة لحمض الساليسيليك:
إنه أسبرين، ولم يُكشف فيه حمض ساليسيليك.
\end{solution}

\begin{solution}{exo:g10:synthesis-yield:10}
\omterm{def:g6:solutions-and-solubility:solute}{المذيب} A: \omterm{def:g7:chemical-reactions:reactant}{الناتج} كثير \omterm{def:g2:mixing-and-dissolving:dissolve}{الذوبان} فيه ساخنًا، ولا يكاد \omterm{def:g2:mixing-and-dissolving:dissolve}{يذوب} فيه باردًا؛ أما \omterm{def:g6:solutions-and-solubility:solute}{المذيب} B
فيُبقي معظم \omterm{def:g7:chemical-reactions:reactant}{الناتج} ذائبًا حتى وهو بارد. ومع A تحتاج \qty{3.0}{g} إلى نحو
$3.0 / 150 = \qty{0.020}{L}$، أي \qty{20}{mL} من \omterm{def:g6:solutions-and-solubility:solute}{المذيب} الساخن؛ وبعد التبريد
يبقى ذائبًا على الأكثر $2 \times 0.020 = \qty{0.04}{g}$.
\end{solution}

\begin{solution}{exo:g10:synthesis-yield:11}
السائل الذي يغلي يحتاج إلى حيز فوقه: فيجب ألا تبلغ الرغوة والرذاذ المكثّف.
وحصى الغليان يعطي البخار أماكن صغيرة تتكوّن فيها الفقاعات، فيغلي السائل بهدوء
بدلًا من أن يغلي في فورات مفاجئة.
\end{solution}

\begin{solution}{exo:g10:synthesis-yield:12}
$0.80 \times 0.70 = 0.56$: \qty{56}{\percent}. وثلاث مراحل \omterm{def:g10:synthesis-yield:yield}{مردود} كل منها
\qty{90}{\percent}: $0.90^3 = 0.73$، أي \qty{73}{\percent}. فكل مرحلة تضيّع
شيئًا من \omterm{def:g7:chemical-reactions:reactant}{الناتج} (ومن الوقت والمال)، والخسائر تتضاعف: وكلما قلّت المراحل كثر
\omterm{def:g7:chemical-reactions:reactant}{الناتج} الباقي في النهاية.
\end{solution}

\begin{solution}{exo:g10:synthesis-yield:13}
على الأكثر $3.3 \times 0.050 = \qty{0.17}{g}$ من الأسبرين. والماء المثلَّج يذيب
من الأسبرين أقل من ذلك، لأن ذوبانيته تنخفض مع
درجة الحرارة.
\end{solution}

\begin{solution}{exo:g10:synthesis-yield:14}
\begin{enumerate}
  \item $n_0(\text{حمض الساليسيليك}) = 2.76 / 138.0 = \qty{0.0200}{mol}$،
        والأنهيدريد \qty{0.0500}{mol}: فحمض الساليسيليك هو المحدِّد،
        $x_{\max} = \qty{0.0200}{mol}$، ويبقى $0.0500 - 0.0200 =
        \qty{0.0300}{mol}$ من الأنهيدريد.
  \item يعطي \omterm{def:g10:synthesis-yield:synthesis}{التخليق} \qty{0.0200}{mol} من حمض الإيثانويك، ويعطي إتلاف
        الزائد $2 \times 0.0300 = \qty{0.0600}{mol}$:
        أي \qty{0.0800}{mol} في المجموع.
\end{enumerate}
\end{solution}

\begin{solution}{exo:g10:synthesis-yield:15}
$n(\text{الأسبرين}) = \num{1.00e6} / 180.0 = \qty{5.56e3}{mol}$. ومع
\omterm{def:g10:synthesis-yield:yield}{مردود} \qty{85}{\percent}، يلزم من حمض الساليسيليك:
$\num{5.56e3} / 0.85 = \qty{6.54e3}{mol}$، أي
$\num{6.54e3} \times 138.0 = \qty{9.0e5}{g}$، أي نحو \qty{0.90}{t}.
\end{solution}

\begin{solution}{pb:g10:synthesis-yield:1}
\textbf{1.} C: $7 + 4 = 11$ و$9 + 2 = 11$؛ H: $6 + 6 = 12$
و$8 + 4 = 12$؛ O: $3 + 3 = 6$ و$4 + 2 = 6$. موزونة.

\textbf{2.} حمض الساليسيليك $7 \times 12.0 + 6 \times 1.0 + 3 \times 16.0
= \qty{138.0}{g/mol}$؛ والأنهيدريد $4 \times 12.0 + 6 \times 1.0 +
3 \times 16.0 = \qty{102.0}{g/mol}$؛ والأسبرين $9 \times 12.0 + 8 \times
1.0 + 4 \times 16.0 = \qty{180.0}{g/mol}$؛ وحمض الإيثانويك $2 \times 12.0
+ 4 \times 1.0 + 2 \times 16.0 = \qty{60.0}{g/mol}$.

\textbf{3.} \omterm{def:g7:chemical-reactions:reactant}{المتفاعلات}: حمض الساليسيليك وأنهيدريد الإيثانويك؛ \omterm{def:g7:chemical-reactions:reactant}{والناتج}
المطلوب: الأسبرين؛ \omterm{def:g7:chemical-reactions:reactant}{والناتج} الثانوي: حمض الإيثانويك.

\textbf{4.} التسخين يسرّع التفاعل؛ وفي الارتداد تُكثَّف الأبخرة وتُعاد، فلا
يضيع أي \omterm{def:g7:chemical-reactions:reactant}{متفاعل} أو \omterm{def:g7:chemical-reactions:reactant}{ناتج}.

\textbf{5.} $n = 3.0 / 138.0 = \qty{0.0217}{mol}$.

\textbf{6.} $m = 6.0 \times 1.08 = \qty{6.48}{g}$؛
و$n = 6.48 / 102.0 = \qty{0.0635}{mol}$.

\textbf{7.} حمض الساليسيليك: $0.0217 - x$؛ والأنهيدريد $0.0635 - x$؛
والأسبرين وحمض الإيثانويك: $x$. ينفد حمض الساليسيليك أولًا:
$x_{\max} = \qty{0.0217}{mol}$، وحمض الساليسيليك هو المحدِّد.

\textbf{8.} $m_{\max} = 0.0217 \times 180.0 = \qty{3.91}{g}$.

\textbf{9.} حتى يتحول حمض الساليسيليك كله إلى أسبرين: فالمتبقي منه يبقى ممزوجًا
بالأسبرين ويصعب نزعه، بينما يُتلف الأنهيدريد الزائد بالماء المضاف في النهاية،
ويُغسل حمض الإيثانويك الذي يعطيه.

\textbf{10.} فهي تشمل الماء: \qty{3.6}{g} تعطي مردودًا زائفًا
$3.6 / 3.91 = \qty{92}{\percent}$.

\textbf{11.} $3.3 / 3.91 = 0.84$: \qty{84}{\percent} من \omterm{def:g10:synthesis-yield:synthesis}{الناتج الخام} الجاف،
وهو غير نقي.

\textbf{12.} تُنزع الشوائب، ويبقى جزء من الأسبرين ذائبًا في \omterm{def:g6:solutions-and-solubility:solute}{المذيب} البارد وفي ماء
الغسل.

\textbf{13.} $3.3 \times 0.030 = \qty{0.10}{g}$ على الأكثر.

\textbf{14.} $\eta = 2.7 / 3.91 = 0.69$: \qty{69}{\percent}.

\textbf{15.} درجة انصهار حادة عند قيمة الأسبرين النقي: فالبلورات أسبرين،
ونقية.

\textbf{16.} انصهار تحت \qty{135}{\celsius} وعلى مدى واسع: \omterm{def:g10:synthesis-yield:synthesis}{فالناتج الخام}
كان غير نقي.

\textbf{17.} كان \omterm{def:g10:synthesis-yield:synthesis}{الناتج الخام} يحوي الأسبرين وحمض ساليسيليك لم يتفاعل؛ وقد نزعت
\omterm{def:g10:synthesis-yield:recrystallisation}{إعادة التبلور} حمض الساليسيليك.

\textbf{18.} $R_f = 3.1 / 5.0 = 0.62$.

\textbf{19.} نحو المقدار نفسه في كل منهما: من \qty{3.91}{g} الممكنة إلى
\qty{3.3}{g} من \omterm{def:g10:synthesis-yield:synthesis}{الناتج الخام}، ومن \qty{3.3}{g} إلى \qty{2.7}{g} في
\omterm{def:g10:synthesis-yield:recrystallisation}{إعادة التبلور}، أي نحو \qty{0.6}{g} في كل منهما. وبما أن \omterm{def:g10:synthesis-yield:synthesis}{الناتج الخام} لم يكن
أسبرين نقيًا، فإن الجزء الأول من العمل أضاع في الحقيقة من الأسبرين أكثر قليلًا
من \qty{0.6}{g}.

\textbf{20.} \omterm{def:g10:synthesis-yield:yield}{مردود} \omterm{def:g10:synthesis-yield:synthesis}{التخليق} نحو \qty{69}{\percent}.
\end{solution}
